% Generated draft: requires mathematical and source review.
\documentclass[11pt]{article}
\usepackage[margin=25mm]{geometry}
\usepackage{fontspec}
\setmainfont{Latin Modern Roman}
\usepackage{amsmath,amssymb,mathtools,unicode-math}
\setmathfont{Latin Modern Math}
\usepackage{enumitem,needspace,xurl,hyperref}
\setlength{\parindent}{0pt}
\setlength{\parskip}{5pt}
\setlength{\emergencystretch}{3em}
\newcommand{\N}{\mathbb{N}}
\newcommand{\Z}{\mathbb{Z}}
\newcommand{\Q}{\mathbb{Q}}
\newcommand{\R}{\mathbb{R}}
\newcommand{\C}{\mathbb{C}}
\newcommand{\F}{\mathbb{F}}
\newcommand{\A}{\mathbb{A}}
\newcommand{\PP}{\mathbb{P}}
\newcommand{\E}{\mathbb{E}}
\newcommand{\eps}{\varepsilon}
\newcommand{\dd}{\,\mathrm{d}}
\begin{document}
\section*{UnsolvedMath ID: 30005503}
\textbf{Source problem code:} catalogue entry
\subsection{Definitions and mathematical statement}
In standard modular representation terminology, a cyclic block is a block with a cyclic defect group: for a prime \(p\), a field \(k\) of characteristic \(p\), and a finite group \(G\), a block algebra has the form \(B=kGe\), where \(e\) is a primitive central idempotent, and its defect groups are required to be cyclic. This terminology does not supply the missing coefficient-field hypotheses or the definition of \(W(B)\) in the record. A finite group \(G\) is quasisimple if \(G=[G,G]\) and \(G/Z(G)\) is a nonabelian simple group.

The source denotes an invariant of a cyclic block \(B\) by \(W(B)\), but the supplied record does not define that invariant, its ambient category, or its equivalence relation. Consequently, the following is a faithful formulation of the requested tasks, conditional on recovering those definitions and the source's modular hypotheses.
\begin{enumerate}
\item For every finite group \(G\) and every cyclic block \(B\) of \(G\), determine whether \(W(B)\) can be recovered from the character table of \(G\). The record does not specify precisely what character-table data are available, how the block is identified in those data, or what constitutes recovery of \(W(B)\).
\item Determine \(W(B)\) for every nonuniserial cyclic block \(B\) of every finite group. The original wording is ``all (non-uniserial) cyclic blocks''; the supplied reconstructed statement takes nonuniserial blocks as the requested class. The role of the parentheses and the precise intended meaning of nonuniserial require clarification.
\item Determine \(W(B)\) for every cyclic block \(B\) of every finite quasisimple group.
\item Reduce the determination requested in item (2) to the quasisimple case requested in item (3). The original text asks whether ``(b)'' can be reduced to ``(c)'', despite numbering its displayed questions (1)--(4). Identifying those references with items (2) and (3) follows the supplied reconstruction and remains a qualification.
\end{enumerate}
No restriction to odd primes is stated for these questions. The odd-prime qualification in the supplied literature assessment concerns reported partial results only. The undefined invariant and the other missing conventions prevent this draft from being a complete mathematical specification.
\subsection{Short English statement}
The problem asks how to recover and classify the invariant \(W(B)\) of cyclic blocks, including nonuniserial blocks and blocks of finite quasisimple groups, and whether the general determination reduces to the quasisimple case. The supplied record does not define \(W(B)\), so the formulation remains conditional and requires mathematical review.
\subsection{Sources}
\begin{itemize}
\item Local UnsolvedMath registry, \texttt{data/unsolvedmath/problems.json}, permanent ID \(30005503\), title ``Source-Algebra Classification of Cyclic Blocks'', exact source problem code \texttt{catalogue entry}. The supplied record attributes the question to the 2023 Oberwolfach workshop report; no individual proposer is supplied.
\item \emph{Representations of Finite Groups} (2023), Oberwolfach Reports, report 19/2023, pp.\ 1031--1084, DOI \texttt{10.4171/owr/2023/19}. \url{https://doi.org/10.4171/owr/2023/19}. This is the authoritative report citation supplied by the record; no exact question page is given.
\item Hiss--Lassueur, ``source algebra equivalence I'', supplied literature-source label. \url{https://www.math.rwth-aachen.de/homes/Gerhard.Hiss/Preprints/sourcealgebraI.pdf}.
\item Hiss--Lassueur, ``source algebra equivalence III'', supplied literature-source label. \url{https://www.math.rwth-aachen.de/~Gerhard.Hiss/Preprints/sourcealgebraIII.pdf}.
\item Supplied dated literature curation, checked \texttt{2026-08-22}: the record labels the problem partially solved and reports that Hiss--Lassueur determine \(W(B)\) from character-table signs for odd primes, classify many quasisimple cases, and have a 2025 continuation treating special-linear and unitary families. It further reports that the four questions remain unsettled in full generality. These are supplied dated assessments, not independently verified mathematical or literature conclusions.
\item The registry labels the statement \texttt{corrected\_verified} and says that it was reconstructed and checked against the official report. That verification claim is supplied provenance; the present draft does not independently verify it or claim to have consulted the linked sources.
\end{itemize}
\end{document}
